\section{Conclusions}

Large-scale citizen science creates a research-data management challenge that extends well beyond the initial collection of observations. In addition to storing participant-generated data, a durable infrastructure must support heterogeneous data types, imperfect records, provenance, privacy protection, subsequent professional enrichment and long-term publication.

ECHOREPO addresses these requirements through an open-source architecture that separates the acquisition environment from the research-data repository and treats ingestion, canonicalisation, quality assurance, recovery, enrichment, privacy protection and FAIR-oriented publication as explicit stages of the data lifecycle.

Its operational deployment with around 9,000 soil samples demonstrates that citizen-generated observations can be integrated with images, laboratory measurements and microbial biodiversity data while preserving a common sample identity. The same deployment also shows that repository-level quality assurance
remains necessary after data collection and illustrates how retained
provenance can be used to investigate and potentially recover problematic
records.

The experience further indicates that publication and interoperability should be designed as part of the research-data workflow rather than added as end-of-project activities. By generating stable machine-readable resources, applying explicit privacy transformations, assigning persistent identifiers and supporting external discovery, ECHOREPO separates the evolving operational repository from research products intended for long-term reuse.

Several broader design principles emerge from this experience: preserve source context sufficiently to support later correction; maintain stable identifiers across the complete sample lifecycle; separate participant-facing acquisition systems from durable research repositories; represent privacy-related transformations explicitly; accommodate progressive enrichment by specialist analyses; and treat persistent publication and external interoperability as architectural concerns.

ECHOREPO remains a soil-specific implementation, and its direct transfer to other scientific domains would require adaptation of the acquisition layer and canonical data model. The broader contribution of this work is therefore not a claim of universal software applicability, but a reusable architectural pattern and a set of operational lessons for transforming heterogeneous citizen-generated observations into quality-controlled, privacy-aware and FAIR-oriented research data.